\documentclass[11pt]{article}
\usepackage[a4paper,margin=27mm]{geometry}
\usepackage{amsmath,amssymb,amsthm,booktabs}
\newtheorem{theorem}{Theorem}
\newtheorem{corollary}[theorem]{Corollary}
\newcommand{\Aut}{\operatorname{Aut}}
\title{The chiral and reflexible pairs of type $\{6,6\}$}
\author{}\date{}
\begin{document}\maketitle
Let $G$ be O'Nan's simple group and $A=\Aut(G)$. We use the convention
of Connor and Leemans, \emph{Algorithmic enumeration of regular maps},
Ars Math. Contemp. 10 (2016), 211--222: generating pairs $(R,S)$ with
$|RS|=2$ are identified under $A$. There is no additional quotient by
interchanging $R$ and $S$, including when their orders agree. An inner
or outer reflexible pair admits an inverter in $G$ or $A\setminus G$,
respectively. The inverter is unique: the ratio of two inverters fixes
the two generators and hence every element of $G$.

\begin{theorem}
There are $165$ classes of type $\{6,6\}$. Of these, $8$ are inner
reflexible, $1$ is outer reflexible, and $156$ are chiral.
\end{theorem}
\begin{proof}
The total $165$ is imported from the independently proved numerical
table. To separate the three kinds, use the established correspondence
between an inverter and a generating triple $(a,b,c)$ of involutions,
where $ac=ca$, $a\ne c$, all three entries have the same parity modulo
$G$, and $R=ab$, $S=bc$. This correspondence is equivariant under $A$.
We reuse its proof and attribution; no new correspondence is asserted.

Since $a,b$ have the same parity, $ab\in G$. If $|ab|=6$, the element
$z=(ab)^3$ is an inner involution and commutes with both $a$ and $b$.
The accepted partition in $C_A(z)$ supplies all four first-pair classes
in the following table. The centralizer of a pair is already contained
in $C_A(z)$, because $z$ is determined by that pair. Thus its acquired
order is the order of the full stabilizer $H=C_A(a,b)$, and no factor
arising from a smaller acting group is missing.
\[
\begin{array}{c|c|r|r|r|r}
\text{class, pair}&\text{parity}&|H|&\text{domain orbits}&M&M/|H|\\\hline
2,1&\text{inner}&4&470&28&7\\
3,1&\text{inner}&16&139&16&1\\
3,2&\text{inner}&16&135&0&0\\
4,1&\text{outer}&12&141&12&1
\end{array}
\]
Here $M$ is the number of eligible third involutions for which
$|bc|=6$ and $\langle ab,bc\rangle=G$.

For completeness, the eligible commuting domains contain $1750$
third involutions when $a$ is inner and $1463$ when $a$ is outer.
The accepted transports put each $a$ at the corresponding fixed
involution. The complete small action of $H$ partitions this domain;
each orbit representative has a literal lift to the acquired action
of degree $245520$. The reader checks $c^2=1$, $c\ne a$, $ac=ca$
and the exact order of $bc$ for every representative. Generation is
checked precisely when that order is six. Representatives of other
orders are explicitly unclassified for generation and do not enter
this calculation.

The generation test is the previously proved orbit criterion. The
socle has a point orbit of size $122760$. None of the thirteen orders
in the exhaustive maximal-subgroup roster is divisible by $122760$.
Consequently a subgroup of $G$ generated by the rotations equals $G$
if and only if its orbit on that point has size $122760$. The reverse
implication follows because a proper subgroup lies in a maximal
subgroup, and each of its point-orbit sizes divides that maximal
subgroup's order. The named faithful model and exhaustive roster are
explicit imported premises.

Every generating triple has trivial stabilizer in $A$, since an
automorphism fixing it fixes the two rotations. In particular each
of its $H$-orbits has size $|H|$; the reader checks this equality for
every generating orbit. The complete partitions give the displayed
values of $M$. Thus the inner count is $7+1+0=8$, and the outer count
is $1$. These are disjoint by uniqueness of the inverter. Subtracting
from the total gives $165-8-1=156$ chiral classes.
\end{proof}

The first independent literal read stopped after $314$ complete orbit
records. Each record was flushed only after its mathematical predicates
passed. The continuation rejoins those representatives in the complete
small partition and checks only the remaining $156$. The two disjoint
intervals are compared with the complete producer journal, as are the
other three independent reads. The incomplete receipt is retained;
it is not presented as a complete acquisition.

\begin{corollary}
Type $\{6,6\}$ admits a chiral map. Together with the previously
accepted $127$ positive chiral types and the absence of chiral maps at
$\{3,15\}$, this proves the source's assertion that exactly $128$
of its $129$ admissible types admit chiral maps.
\end{corollary}

The exact split in the theorem refines this existence assertion; it is
not read from the source's claimed number of types. It also supplies
an inner and an outer witness of type $\{6,6\}$. Existence of all
$95$ inner types and all $98$ outer types remains a separate spectrum
question. This component alone does not grant whole-source credit.
\end{document}
